\documentclass[11pt]{article}
%=======================================
% NORTHWESTERN -- Personal Statement (125--250 words)
% Prompt: "Write a statement that includes any relevant information on one or
% more of the following topics (the length of this statement should range
% from 125 to 250 words maximum):
% - Pivotal experiences, opportunities, and/or challenges that have influenced
%   your educational and professional development.
% - Leadership experiences, community outreach, service initiatives, or
%   research projects you have participated in or plan to pursue that aim to
%   positively impact others or the broader community.
% - Anything you would like to share with the admissions committee that you
%   have yet to discuss in other areas of the application can be addressed
%   here. This may include any gaps in your academic and professional
%   experience, or additional context if you believe your academic record
%   does not demonstrate your true capabilities. If needed, you may also use
%   this space to expand upon the topics discussed in the above essays."
%
% Any one topic is enough, so this uses the generic background pieces only
% (no leadership section).
% The shared pieces are well over 250 words, so this will need cut-down
% versions (e.g. master/01_background_short.tex).
%=======================================
\newcommand{\school}{Northwestern University}
\newcommand{\wordmin}{125}
\newcommand{\wordlimit}{250}
%TC:macro \placeholder [ignore]   % keep red notes out of the word count
%TC:ignore
%=======================================
% Shared preamble for every personal statement.
% Each program file sets \school, \pstitle and (optionally) \wordmin and
% \wordlimit BEFORE \input-ing this file; anything left unset falls back to
% the defaults below.
%=======================================
\usepackage{lmodern}          % AER font
\usepackage[T1]{fontenc}      % Better font encoding
\usepackage{amsmath, amssymb}
\usepackage[margin=1in]{geometry}
\usepackage{graphicx}
\usepackage{url}
\usepackage{xcolor}
\usepackage[round]{natbib}
\usepackage{bibentry}         % lets \nobibliography resolve cites without printing a list
\usepackage{hyperref}
\usepackage{fancyhdr}

%=======================================
% Format helpers
%=======================================
\providecommand{\school}{[University name]}
\providecommand{\program}{Economics}
\providecommand{\pstitle}{Personal Statement}
\providecommand{\schoolheader}{\school}   % header line; defaults to \school
% Path from the compiling file's folder back to the repo root (school files
% sit in schools/<school>/tex/; Master_PS overrides this).
\providecommand{\reporoot}{../../../}

% Run-in section heading: bold label ending in a period, body text follows
% on the same line (matches the reference format).
\newcommand{\shead}[1]{\par\vspace{0.85em}\noindent\textbf{#1.}\ }

% Suppress the indent of whichever paragraph starts next. Unlike \noindent
% this survives blank lines/comments in between, and does nothing if no
% paragraph follows (so a trailing \psbreak can't create an empty page).
\newcommand{\psnoindent}{\everypar{{\setbox0\lastbox}\everypar{}}}

% Break between master pieces: a blank line, next line not indented.
\newcommand{\psbreak}{\par\vspace{\baselineskip}\psnoindent}

% Red bracketed note for text still to be written. Each program file tells
% texcount to skip it (see the %TC:macro line in each program file).
\newcommand{\placeholder}[1]{\textcolor{red}{[#1]}}

% Indented paragraphs, no extra vertical space between them.
\setlength{\parindent}{1.5em}
\setlength{\parskip}{0pt}

%=======================================
% Footer: page number centred, word count on the right.
% The count is read from <jobname>.wc, which the texcount build recipe
% regenerates on each build. Text between %TC:ignore and %TC:endignore is
% excluded from the count. The count turns red when outside the limits.
%=======================================
\newread\wcfile
\newif\ifwcknown
\IfFileExists{\jobname.wc}{%
  \openin\wcfile=\jobname.wc
  \read\wcfile to\wcnum
  \closein\wcfile
  \wcknowntrue}{\def\wcnum{??}}

\newif\ifwcbad
\newcommand{\wordcountdisplay}{%
  \wcbadfalse
  \ifwcknown
    \ifdefined\wordlimit\ifnum\wcnum>\wordlimit\relax\wcbadtrue\fi\fi
    \ifdefined\wordmin\ifnum\wcnum<\wordmin\relax\wcbadtrue\fi\fi
  \fi
  \ifwcbad\color{red}\fi
  Word count: \wcnum%
  \ifdefined\wordlimit
  %
  \fi}

\pagestyle{fancy}
\fancyhf{}
\renewcommand{\headrulewidth}{0pt}
\fancyfoot[C]{\thepage}
\fancyfoot[R]{\small\wordcountdisplay}
%TC:endignore


\begin{document}
%TC:ignore
%=======================================
% Header (not counted)
%=======================================
\noindent
\begin{minipage}[t]{0.60\textwidth}
  {\large\bfseries \pstitle}
\end{minipage}\hfill
\begin{minipage}[t]{0.38\textwidth}
  \raggedleft
  Rhea Choudhary\\[2pt]
  \schoolheader
\end{minipage}

\vspace{0.5em}
\noindent\rule{\textwidth}{0.8pt}

\vspace{1.1em}
\psnoindent
%TC:endignore


% Northwestern-specific text (250-word max): the "As I look back" paragraph,
% a one-sentence version of where support falls short, then why grad school.
% Edits to the shared master pieces do NOT flow through here.
As I look back on my high school and university experiences, I think a lot about how important my family and the informal safety net they provided was to my journey to graduate school. I could take risks in my career, take on more study or choose more demanding jobs because I had informal support. But I also think a lot about how getting through the door to these opportunities required government support; I went to a public university, used government-issued university loans, and relied on income assistance when studying full-time. Both government assistance and family support played complementary roles in helping my career.
\psbreak
But I've also seen where both fall short: my brother was unemployed for over three years and needed more than the imperfect income assistance provided by unemployment insurance, and government programs aren't well-positioned to alleviate labour-market scarring in the face of technological disruption.
\psbreak
These experiences have strongly motivated my decision to pursue graduate study, because I want to answer important questions like these that affect people's lives, and are front-of-mind within the academic community as well as in policymaking. I want to expand my horizons and eventually push the frontier of the discipline to answer questions of deep importance to me.

%TC:ignore
% Citations resolve against bib_PS.bib, but no reference list is printed.
% \nocite{*} keeps bibtex from erroring when a statement cites nothing.
\nocite{*}
\bibliographystyle{abbrvnat}
\nobibliography{\reporoot shared/ps/bib_PS}
%TC:endignore

\end{document}
